pdfoutput=1
% Options for packages loaded elsewhere
\PassOptionsToPackage{unicode}{hyperref}
\PassOptionsToPackage{hyphens}{url}
\pdfoutput=1
\documentclass[
  12pt,
]{article}
\usepackage{xcolor}
\usepackage{amsmath,amssymb}
\setcounter{secnumdepth}{-\maxdimen} % remove section numbering
\usepackage{iftex}
\ifPDFTeX
  \usepackage[T1]{fontenc}
  \usepackage{textcomp} % provide euro and other symbols
\else % if luatex or xetex
  \usepackage{unicode-math} % this also loads fontspec
  \defaultfontfeatures{Scale=MatchLowercase}
  \defaultfontfeatures[\rmfamily]{Ligatures=TeX,Scale=1}
\fi
\usepackage{lmodern}
\ifPDFTeX\else
  % xetex/luatex font selection
\fi
% Use upquote if available, for straight quotes in verbatim environments
\IfFileExists{upquote.sty}{\usepackage{upquote}}{}
\IfFileExists{microtype.sty}{% use microtype if available
  \usepackage[]{microtype}
  \UseMicrotypeSet[protrusion]{basicmath} % disable protrusion for tt fonts
}{}
\makeatletter
\@ifundefined{KOMAClassName}{% if non-KOMA class
  \IfFileExists{parskip.sty}{%
    \usepackage{parskip}
  }{% else
    \setlength{\parindent}{0pt}
    \setlength{\parskip}{6pt plus 2pt minus 1pt}}
}{% if KOMA class
  \KOMAoptions{parskip=half}}
\makeatother
\usepackage{longtable,booktabs,array}
\usepackage{calc} % for calculating minipage widths
% Correct order of tables after \paragraph or \subparagraph
\usepackage{etoolbox}
\makeatletter
\patchcmd\longtable{\par}{\if@noskipsec\mbox{}\fi\par}{}{}
\makeatother
% Allow footnotes in longtable head/foot
\IfFileExists{footnotehyper.sty}{\usepackage{footnotehyper}}{\usepackage{footnote}}
\makesavenoteenv{longtable}
\setlength{\emergencystretch}{3em} % prevent overfull lines
\providecommand{\tightlist}{%
  \setlength{\itemsep}{0pt}\setlength{\parskip}{0pt}}
\usepackage{bookmark}
\IfFileExists{xurl.sty}{\usepackage{xurl}}{} % add URL line breaks if available
\urlstyle{same}
\hypersetup{
  hidelinks,
  pdfcreator={LaTeX via pandoc}}

\author{SCX}
\date{2026}

\begin{document}

\section{SCX Self-Evolution Theory: Verification
Report}\label{scx-self-evolution-theory-verification-report}

\begin{quote}
\textbf{Part of the SCX Self-Evolution Theory Series} \textbf{Status}:
Pre-review verification \textbf{Version}: 2026-06-28 \textbf{Purpose}:
Internal consistency check, cross-reference with existing theory, gap
analysis, and honest assessment
\end{quote}

\begin{center}\rule{0.5\linewidth}{0.5pt}\end{center}

\subsection{Table of Contents}\label{table-of-contents}

\begin{enumerate}
\def\labelenumi{\arabic{enumi}.}
\tightlist
\item
  \hyperref[1-consistency-check-of-new-definitions]{Consistency Check of
  New Definitions}
\item
  \hyperref[2-cross-reference-with-existing-theorems-1-5]{Cross-Reference
  with Existing Theorems 1-5}
\item
  \hyperref[3-notation-conflict-check]{Notation Conflict Check}
\item
  \hyperref[4-assumption-dependency-map]{Assumption Dependency Map}
\item
  \hyperref[5-proof-gaps-identified]{Proof Gaps Identified}
\item
  \hyperref[6-open-problems]{Open Problems}
\item
  \hyperref[7-numerical-verification-suggestions]{Numerical Verification
  Suggestions}
\item
  \hyperref[8-honest-assessment]{Honest Assessment}
\item
  \hyperref[9-summary-of-verification-status]{Summary of Verification
  Status}
\end{enumerate}

\begin{center}\rule{0.5\linewidth}{0.5pt}\end{center}

\subsection{1. Consistency Check of New
Definitions}\label{consistency-check-of-new-definitions}

\subsubsection{1.1 Definition Inventory}\label{definition-inventory}

The self-evolution theory introduces the following new definitions:

\begin{longtable}[]{@{}
  >{\raggedright\arraybackslash}p{(\linewidth - 6\tabcolsep) * \real{0.2353}}
  >{\raggedright\arraybackslash}p{(\linewidth - 6\tabcolsep) * \real{0.3235}}
  >{\raggedright\arraybackslash}p{(\linewidth - 6\tabcolsep) * \real{0.1765}}
  >{\raggedright\arraybackslash}p{(\linewidth - 6\tabcolsep) * \real{0.2647}}@{}}
\toprule\noalign{}
\begin{minipage}[b]{\linewidth}\raggedright
Symbol
\end{minipage} & \begin{minipage}[b]{\linewidth}\raggedright
Definition
\end{minipage} & \begin{minipage}[b]{\linewidth}\raggedright
Type
\end{minipage} & \begin{minipage}[b]{\linewidth}\raggedright
Used In
\end{minipage} \\
\midrule\noalign{}
\endhead
\bottomrule\noalign{}
\endlastfoot
\(S_t\) & Gatekeeper scoring function at time \(t\) & Function
\(\mathcal{X} \to [0,1]\) & SE-1, SE-2, SE-C1 \\
\(M_t\) & Memory bank at time \(t\) & Set of data triples & SE-1,
SE-2 \\
\(\theta_t\) & NEP student model parameters at time \(t\) & Parameter
vector & SE-1, SE-2 \\
\(\Phi(S_t, M_t, f_{\theta_t})\) & Lyapunov function for convergence &
Scalar & SE-1, SE-2 \\
\(\rho_s\) & State probability (carried over from Thm 1) & Scalar
\(\in [0,1]\) & SE-1 \\
\(V(s)\) & State data value (carried over from definitions/04) & Scalar
& SE-1, connections file \\
\(\mathcal{F}\) & Gatekeeper function class & Set of functions & SE-2,
SE-4, SE-5 \\
\(\mathcal{F}_{\text{dist}}\) & Distinguishable gatekeeper functions &
Finite set & SE-3 \\
\(\mathcal{Q}\) & System configuration space & Finite set & SE-2,
SE-3 \\
\(\varepsilon_{\text{mach}}\) & Machine precision & Scalar \(> 0\) &
SE-2, SE-3 \\
\(\gamma_t\) & Lyapunov descent step & Scalar \(\geq 0\) & SE-2 \\
\(\mathcal{A}\) & Action space for state-certified AL & Set of actions &
File 08 \\
\end{longtable}

\subsubsection{1.2 Internal Consistency
Checks}\label{internal-consistency-checks}

\textbf{Check 1: \(S_t\) domain and range.} - Domain: \(\mathcal{X}\)
(input space) -- consistent with existing theory (\(\mathcal{X}\) is the
input space across all theorems). - Range: \([0,1]\) -- consistent with
being a scoring function that estimates label correctness probability. -
Dependency: \(S_t\) depends on \(M_t\) (memory bank) and the state
partition from clustering, not directly on \(\theta_t\) (though both
share the partition). - \textbf{Verdict}: Consistent.

\textbf{Check 2: \(M_t\) monotonicity.} - By definition,
\(M_t \subseteq M_{t+1}\): data is only added, never removed. - Bounded
above by \(|\mathcal{D}_{\text{total}}| \leq N_{\text{max}}\). - This
monotonicity is critical for Theorem SE-2's convergence argument. -
\textbf{Verdict}: Consistent.

\textbf{Check 3: Lyapunov function \(\Phi\) domain.} - \(\Phi\) maps
\((S_t, M_t, f_{\theta_t})\) to \(\mathbb{R}_{\geq 0}\). - The three
arguments are all well-defined objects. The value depends on: -
Gatekeeper quality (measured by consensus-reliability alignment) -
Memory bank quality (measured by coverage and diversity) - Student
quality (measured by NEP prediction error on validated data) -
\textbf{Verdict}: The function is well-defined but its exact form is
\textbf{underspecified} in the current files. The Lyapunov function must
be concretely defined for the theory to be fully rigorous. See Section
5.1.

\textbf{Check 4: \(\varepsilon_{\text{mach}}\) usage.} - Used in two
contexts: (a) as machine precision for numerical distinguishability; (b)
as the minimum Lyapunov descent step. - These are related: if numerical
values cannot be distinguished below \(\varepsilon_{\text{mach}}\), then
\(\gamma_t < \varepsilon_{\text{mach}}\) is indistinguishable from
\(\gamma_t = 0\). - \textbf{Verdict}: Consistent across uses.

\textbf{Check 5: State space \(\mathcal{S}\) vs.~\(\mathcal{S}_t\).} -
Existing theory: \(\mathcal{S}\) is the fixed, true state space
(partition of \(\mathcal{X}\) into \(K_S\) states). - Self-evolution:
The state partition may be refined or updated as the gatekeeper
improves. - The current self-evolution files use \(\mathcal{S}\) to
denote the current best estimate of the partition, which may evolve with
\(t\). - \textbf{Potential conflict}: If \(\mathcal{S}\) denotes the
true (unknown) state space, it should not be time-indexed. If it denotes
the estimated partition, it should be \(\mathcal{S}_t\) or
\(\hat{\mathcal{S}}_t\). - \textbf{Recommendation}: Use
\(\hat{\mathcal{S}}_t\) for the estimated state partition at time \(t\),
reserving \(\mathcal{S}\) for the true partition. This recommendation
applies to all self-evolution files. The current files use
\(\mathcal{S}\) ambiguously.

\subsubsection{1.3 Consistency Summary}\label{consistency-summary}

\begin{longtable}[]{@{}
  >{\raggedright\arraybackslash}p{(\linewidth - 4\tabcolsep) * \real{0.3182}}
  >{\raggedright\arraybackslash}p{(\linewidth - 4\tabcolsep) * \real{0.3636}}
  >{\raggedright\arraybackslash}p{(\linewidth - 4\tabcolsep) * \real{0.3182}}@{}}
\toprule\noalign{}
\begin{minipage}[b]{\linewidth}\raggedright
Check
\end{minipage} & \begin{minipage}[b]{\linewidth}\raggedright
Status
\end{minipage} & \begin{minipage}[b]{\linewidth}\raggedright
Notes
\end{minipage} \\
\midrule\noalign{}
\endhead
\bottomrule\noalign{}
\endlastfoot
\(S_t\) domain/range & PASS & Consistent with
\(\mathcal{X} \to [0,1]\) \\
\(M_t\) monotonicity & PASS & Well-defined and bounded \\
\(\Phi\) exact form & \textbf{GAP} & Needs explicit definition \\
\(\varepsilon_{\text{mach}}\) usage & PASS & Consistently applied \\
\(\mathcal{S}\) vs \(\mathcal{S}_t\) & \textbf{WARNING} & Ambiguous;
recommend \(\hat{\mathcal{S}}_t\) for estimated partition \\
\end{longtable}

\begin{center}\rule{0.5\linewidth}{0.5pt}\end{center}

\subsection{2. Cross-Reference with Existing Theorems
1-5}\label{cross-reference-with-existing-theorems-1-5}

\subsubsection{2.1 Theorem SE-1 vs.~Theorem 1 (Noise
Detection)}\label{theorem-se-1-vs.-theorem-1-noise-detection}

\textbf{Question}: Does Theorem SE-1 rely on or contradict Theorem 1?

\textbf{Analysis}: - Theorem SE-1 uses the consensus score
\(C(x) = \frac{1}{M}\sum_m \mathbf{1}\{\ell(f_m(x), y) > \tau\}\) (same
definition as Theorem 1). - Theorem SE-1's Lyapunov function \(\Phi\)
should include a term that measures the gap between the gatekeeper's
score \(S_t(x)\) and the consensus-based noise estimate \(C(x)\). - When
Theorem 1's conditions hold (\(\Delta_s > 0\) for all states), the
gatekeeper can reliably use \(C(x)\) as a signal, and the Lyapunov
descent is well-conditioned. - When Theorem 1 fails (\(\Delta_s \leq 0\)
for some states), the consensus signal is unreliable in those states,
and the Lyapunov gradient may point in the wrong direction.

\textbf{Conclusion}: Theorem SE-1 does \textbf{not contradict} Theorem
1: - Theorem 1 provides the \textbf{signal strength} that powers the
Lyapunov descent. - Theorem SE-1's convergence \textbf{depends on}
Theorem 1's condition: if \(\Delta_s = 0\) for all states, the Lyapunov
gradient vanishes and convergence stalls. - This dependency should be
made explicit: Theorem SE-1 should state its result conditionally on the
existence of states with \(\Delta_s > 0\).

\textbf{Verdict}: Consistent (with explicit dependency needed).

\subsubsection{2.2 Self-Evolution vs.~Theorem 2 (Weak Feature
Failure)}\label{self-evolution-vs.-theorem-2-weak-feature-failure}

\textbf{Question}: Does the self-evolution framework respect Theorem 2's
limitation? Can it overcome weak features through iteration?

\textbf{Analysis}: - Theorem 2 states:
\(F1_{\text{SCX}} \leq F1_{\text{base}} + C_F\sqrt{\delta/2}\) where
\(\delta = I(\phi; S)\) is the feature-state mutual information. - This
is a \textbf{static} bound that depends only on the feature
representation \(\phi\), not on the number of evolution steps \(t\). -
The self-evolution loop \textbf{does not change \(\phi\)}: the feature
representation is fixed throughout the evolution. (The NEP student
\(f_{\theta_t}\) may learn better internal representations, but the
clustering and gatekeeper operate on \(\phi\).) - Therefore, Theorem 2's
bound applies equally to every iteration \(t\):
\[F1_{\text{SCX}, t} \leq F1_{\text{base}} + C_F\sqrt{\delta/2}, \quad \forall t \geq 0\]

\textbf{Conclusion}: The self-evolution framework \textbf{respects}
Theorem 2's limitation. No amount of iteration can overcome the feature
bottleneck. If \(\delta \approx 0\), the fixed point \(S_{T^*}\) cannot
outperform the loss baseline.

\textbf{Implication}: The self-evolution framework should include a
``feature escape'' mechanism in future work, where \(\phi\) is updated
based on the NEP student's learned representations. Without this, the
system is fundamentally bottlenecked by the initial feature choice.

\textbf{Verdict}: Consistent (with feature escape as future work).

\subsubsection{2.3 Self-Evolution vs.~Theorem 3
(Unidentifiability)}\label{self-evolution-vs.-theorem-3-unidentifiability}

\textbf{Question}: Does the self-evolution framework progressively
resolve Theorem 3's unidentifiability?

\textbf{Analysis}: - Theorem 3 states that without assumptions A1-A6,
noise and difficulty are indistinguishable. Even with A1-A6, there
exists an irreducible ambiguity proportional to \(\eta\rho/2\). - The
self-evolution framework \textbf{does not add new assumptions} that
would break the unidentifiability. The same ambiguity persists. -
However, the self-evolution loop may \textbf{reduce the effective
ambiguity} through iteration: - Initial iterations may flag samples as
noise or difficulty based on limited evidence. - As the memory bank
grows, the gatekeeper's estimates become more accurate (Theorem 1's
bound improves with \(M_t\)'s growing evidence). - The irreducible
ambiguity \(\eta\rho/2\) remains, but the \textbf{direction of error}
(false positives vs.~false negatives) may shift as the gatekeeper
calibrates.

\textbf{Claim}: The self-evolution framework \textbf{partially resolves}
the practical impact of unidentifiability by: 1. Explicitly tracking
uncertainty via the consensus score \(C(x)\). 2. Selecting samples for
NEP validation at the intersection of high-uncertainty and high-impact
states. 3. Accumulating validated labels in \(M_t\), which provides
ground truth for recalibration.

\textbf{Limitation}: The self-evolution framework can only resolve
ambiguity that \textbf{data can resolve}. The fundamental \(\eta\rho/2\)
lower bound from Theorem 3 remains: there will always be samples for
which noise vs.~difficulty is undecidable given the available evidence.

\textbf{Verdict}: Consistent. Self-evolution reduces practical ambiguity
but cannot eliminate the fundamental bound.

\subsubsection{2.4 Self-Evolution vs.~Theorem 4' (Exact
Constant)}\label{self-evolution-vs.-theorem-4-exact-constant}

\textbf{Question}: Does the self-evolution framework align with Theorem
4'\,'s optimality claims?

\textbf{Analysis}: - Theorem 4' shows that SCX with an adaptive
threshold achieves the exact constant minimax optimal rate for noise
detection. - The self-evolution framework's gatekeeper \(S_t\) is a
generalization of the fixed threshold detector analyzed in Theorem 4'. -
At the fixed point \(T^*\), the gatekeeper \(S_{T^*}\) should
approximate the optimal threshold \(\theta^\dagger\) from Theorem 4'. -
The Lyapunov convergence (Theorem SE-1) should drive \(S_t\) toward
\(\theta^\dagger\) in states where the expert consensus is
well-conditioned.

\textbf{Connection}: Theorem 4'\,'s \(C_{\min}/\eta\) optimal constant
provides the \textbf{target} for the self-evolution's fixed point. If
the Lyapunov function is well-designed, the fixed point should achieve
approximately this optimal constant.

\textbf{Implication}: The self-evolution fixed point's quality is
bounded by Theorem 4'\,'s optimality constant. This provides a
\textbf{stopping criterion}: when \(S_t(x) \approx \theta^\dagger\) for
all \(x\) with \(\Delta_s > 0\), further evolution cannot significantly
improve noise detection.

\textbf{Verdict}: Consistent. The connection should be made explicit in
the definition of \(\Phi\).

\subsubsection{2.5 Self-Evolution vs.~Theorem 5 and Proposition
6}\label{self-evolution-vs.-theorem-5-and-proposition-6}

\textbf{Question}: How does self-evolution depend on state discovery
quality?

\textbf{Analysis}: - Theorem 5 guarantees that k-means recovers the true
state partition under strong separation. - Proposition 6 provides a
bootstrap diagnostic for whether the state structure is meaningful. -
The self-evolution framework relies on the state partition for the
gatekeeper's state-level statistics. - If Theorem 5's conditions hold
(\(\Delta_{\min}\) sufficiently large), the gatekeeper operates on a
correct partition and its estimates are accurate. - If Proposition 6's
stability score is low, the gatekeeper's state-level statistics are
unreliable, and the Lyapunov descent may be misdirected.

\textbf{Recommendation}: The self-evolution loop should check
Proposition 6's stability criterion as a \textbf{precondition} before
relying on state-level statistics. This is currently not formalized in
the self-evolution files.

\textbf{Verdict}: Consistent (with precondition check recommended).

\subsubsection{2.6 Cross-Reference
Summary}\label{cross-reference-summary}

\begin{longtable}[]{@{}
  >{\raggedright\arraybackslash}p{(\linewidth - 4\tabcolsep) * \real{0.3269}}
  >{\raggedright\arraybackslash}p{(\linewidth - 4\tabcolsep) * \real{0.5192}}
  >{\raggedright\arraybackslash}p{(\linewidth - 4\tabcolsep) * \real{0.1538}}@{}}
\toprule\noalign{}
\begin{minipage}[b]{\linewidth}\raggedright
Existing Theorem
\end{minipage} & \begin{minipage}[b]{\linewidth}\raggedright
Relation to Self-Evolution
\end{minipage} & \begin{minipage}[b]{\linewidth}\raggedright
Status
\end{minipage} \\
\midrule\noalign{}
\endhead
\bottomrule\noalign{}
\endlastfoot
Thm 1 (Noise Detection) & Provides gradient signal for Lyapunov descent
& Consistent (dependency should be explicit) \\
Thm 2 (Weak Feature) & Fundamental bound that iteration cannot overcome
& Consistent (feature escape needed for future) \\
Thm 3 (Unidentifiability) & Irreducible ambiguity persists; iteration
reduces practical impact & Consistent \\
Thm 4' (Exact Constant) & Fixed point target; optimality constant bounds
final quality & Consistent (connection should be explicit) \\
Thm 5 (Cluster Consistency) & State discovery quality bounds gatekeeper
accuracy & Consistent (precondition check recommended) \\
Prop 6 (Stability Diagnostic) & Precondition for reliable state-level
statistics & Consistent (precondition check recommended) \\
\end{longtable}

\begin{center}\rule{0.5\linewidth}{0.5pt}\end{center}

\subsection{3. Notation Conflict Check}\label{notation-conflict-check}

\subsubsection{3.1 Symbol Usage Across Self-Evolution and Existing
Theory}\label{symbol-usage-across-self-evolution-and-existing-theory}

\begin{longtable}[]{@{}
  >{\raggedright\arraybackslash}p{(\linewidth - 6\tabcolsep) * \real{0.1379}}
  >{\raggedright\arraybackslash}p{(\linewidth - 6\tabcolsep) * \real{0.2931}}
  >{\raggedright\arraybackslash}p{(\linewidth - 6\tabcolsep) * \real{0.3793}}
  >{\raggedright\arraybackslash}p{(\linewidth - 6\tabcolsep) * \real{0.1897}}@{}}
\toprule\noalign{}
\begin{minipage}[b]{\linewidth}\raggedright
Symbol
\end{minipage} & \begin{minipage}[b]{\linewidth}\raggedright
Existing Meaning
\end{minipage} & \begin{minipage}[b]{\linewidth}\raggedright
Self-Evolution Meaning
\end{minipage} & \begin{minipage}[b]{\linewidth}\raggedright
Conflict?
\end{minipage} \\
\midrule\noalign{}
\endhead
\bottomrule\noalign{}
\endlastfoot
\(\mathcal{S}\) & True state space (fixed) & Current estimated state
space (may vary with \(t\)) & \textbf{YES} -- see recommendation in
Section 1.2, Check 5 \\
\(S_t\) & Not used & Gatekeeper scoring function & None (new symbol) \\
\(M\) & Number of experts & Same & None \\
\(M_t\) & Not used & Memory bank at time \(t\) & None (new symbol) \\
\(\theta\) & Detection threshold (Thm 1,4') & NEP student parameters
\(\theta_t\) & \textbf{YES} -- \(\theta\) is overloaded \\
\(\Phi\) & Feature representation \(\phi(X)\) (Thm 2,5) & Lyapunov
function & \textbf{YES} -- \(\Phi\) (capital phi) vs \(\phi\)
(lowercase) are different but easily confused \\
\(\rho_s\) & State probability (Thm 1) & Same & None \\
\(f_\theta\) & Not used in existing theory & NEP student model & None
(new) \\
\(\mathcal{F}\) & Not used explicitly & Gatekeeper function class & None
(new) \\
\(\eta\) & Global noise rate (Thm 1,2,3,4') & Same & None \\
\(\tau\) & Expert error threshold (Thm 1) & Same & None \\
\end{longtable}

\subsubsection{3.2 Resolution of
Conflicts}\label{resolution-of-conflicts}

\textbf{Conflict 1: \(\mathcal{S}\) vs.~\(\mathcal{S}_t\)}: -
\textbf{Recommendation}: Use \(\hat{\mathcal{S}}_t\) (estimated state
partition at time \(t\)) in self-evolution files. Reserve
\(\mathcal{S}\) for the true partition. - \textbf{Action needed}: Update
all self-evolution files to use \(\hat{\mathcal{S}}_t\) where the
estimated partition is meant. Files 07 and 08 currently use
\(\mathcal{S}\) without time indexing. This is load-bearing for Theorem
SE-2's configuration space finiteness argument. - \textbf{Priority}:
High. This affects the rigor of the configuration space definition.

\textbf{Conflict 2: \(\theta\) overloaded}: - Existing: \(\theta\) =
detection threshold (Theorem 1, 4'). - Self-evolution: \(\theta_t\) =
NEP student model parameters. - \textbf{Resolution}: These are distinct.
The student parameters are naturally distinguished by the subscript
\(t\) and by context (used in \(f_{\theta_t}\), not as a scalar
threshold). No change needed, but users should be aware of the
overloading. - \textbf{Priority}: Low. Context disambiguates.

\textbf{Conflict 3: \(\Phi\) vs.~\(\phi\)}: - Existing: \(\phi(X)\) =
feature representation (lowercase phi). - Self-evolution:
\(\Phi(S_t, M_t, f_{\theta_t})\) = Lyapunov function (uppercase Phi). -
\textbf{Resolution}: These are distinct Greek letters (\(\Phi\) vs
\(\phi\)) and used in different contexts. However, the proximity could
cause confusion in formula-dense text. - \textbf{Recommendation}: The
self-evolution files should always write \(\Phi(\cdot)\) with explicit
arguments to avoid confusion with \(\phi(X)\). - \textbf{Priority}: Low.
Mostly a readability concern.

\subsubsection{3.3 Symbol Collision Risk
Summary}\label{symbol-collision-risk-summary}

\begin{longtable}[]{@{}
  >{\raggedright\arraybackslash}p{(\linewidth - 6\tabcolsep) * \real{0.2391}}
  >{\raggedright\arraybackslash}p{(\linewidth - 6\tabcolsep) * \real{0.1739}}
  >{\raggedright\arraybackslash}p{(\linewidth - 6\tabcolsep) * \real{0.2174}}
  >{\raggedright\arraybackslash}p{(\linewidth - 6\tabcolsep) * \real{0.3696}}@{}}
\toprule\noalign{}
\begin{minipage}[b]{\linewidth}\raggedright
Risk Level
\end{minipage} & \begin{minipage}[b]{\linewidth}\raggedright
Symbol
\end{minipage} & \begin{minipage}[b]{\linewidth}\raggedright
Conflict
\end{minipage} & \begin{minipage}[b]{\linewidth}\raggedright
Recommended Fix
\end{minipage} \\
\midrule\noalign{}
\endhead
\bottomrule\noalign{}
\endlastfoot
\textbf{HIGH} & \(\mathcal{S}\) & True partition vs.~estimated partition
& \(\hat{\mathcal{S}}_t\) for estimated \\
\textbf{LOW} & \(\theta\) & Threshold vs.~student params & Contextual
(acceptable) \\
\textbf{LOW} & \(\Phi\) vs \(\phi\) & Lyapunov vs.~features & Explicit
arguments (acceptable) \\
\textbf{NONE} & All others & No conflicts & -- \\
\end{longtable}

\begin{center}\rule{0.5\linewidth}{0.5pt}\end{center}

\subsection{4. Assumption Dependency
Map}\label{assumption-dependency-map}

\subsubsection{4.1 Existing Assumptions (A1-A6) Needed for
Self-Evolution}\label{existing-assumptions-a1-a6-needed-for-self-evolution}

\begin{longtable}[]{@{}
  >{\raggedright\arraybackslash}p{(\linewidth - 4\tabcolsep) * \real{0.4400}}
  >{\raggedright\arraybackslash}p{(\linewidth - 4\tabcolsep) * \real{0.3600}}
  >{\raggedright\arraybackslash}p{(\linewidth - 4\tabcolsep) * \real{0.2000}}@{}}
\toprule\noalign{}
\begin{minipage}[b]{\linewidth}\raggedright
Assumption
\end{minipage} & \begin{minipage}[b]{\linewidth}\raggedright
Needed?
\end{minipage} & \begin{minipage}[b]{\linewidth}\raggedright
Why
\end{minipage} \\
\midrule\noalign{}
\endhead
\bottomrule\noalign{}
\endlastfoot
\textbf{A1}: Disjoint training sets & Partial & Needed if the experts
\(\{f_m\}\) are used for consensus scoring in the gatekeeper. Not needed
if the gatekeeper uses only the NEP student predictions. \\
\textbf{A2}: Conditional independence (clean) & Partial & Same as A1.
Needed for consensus-based detection; not needed for NEP student
training. \\
\textbf{A3}: Bounded loss & Yes & Required for concentration bounds in
Theorem 1, which powers the Lyapunov gradient. Also needed for NEP
student loss bounds. \\
\textbf{A4}: Uniform independent noise & Partial & Needed for the
noise-detection interpretation of consensus scores. Not strictly needed
if the gatekeeper uses other signals. \\
\textbf{A5}: State homogeneity & Yes & Required for state-level
statistics to be meaningful. Without it, within-state variance makes the
gatekeeper's state-level decisions unreliable. \\
\textbf{A6}: Balanced error distribution & Partial & Needed for the
noise-side concentration bound (Theorem 1, Lemma 3). The self-evolution
framework can operate without it but with degraded guarantees. \\
\end{longtable}

\subsubsection{4.2 New Assumptions Introduced by
Self-Evolution}\label{new-assumptions-introduced-by-self-evolution}

The self-evolution theory introduces the following new assumptions:

\begin{longtable}[]{@{}
  >{\raggedright\arraybackslash}p{(\linewidth - 8\tabcolsep) * \real{0.0638}}
  >{\raggedright\arraybackslash}p{(\linewidth - 8\tabcolsep) * \real{0.2340}}
  >{\raggedright\arraybackslash}p{(\linewidth - 8\tabcolsep) * \real{0.3617}}
  >{\raggedright\arraybackslash}p{(\linewidth - 8\tabcolsep) * \real{0.1915}}
  >{\raggedright\arraybackslash}p{(\linewidth - 8\tabcolsep) * \real{0.1489}}@{}}
\toprule\noalign{}
\begin{minipage}[b]{\linewidth}\raggedright
\#
\end{minipage} & \begin{minipage}[b]{\linewidth}\raggedright
Assumption
\end{minipage} & \begin{minipage}[b]{\linewidth}\raggedright
Formal Statement
\end{minipage} & \begin{minipage}[b]{\linewidth}\raggedright
Used In
\end{minipage} & \begin{minipage}[b]{\linewidth}\raggedright
Notes
\end{minipage} \\
\midrule\noalign{}
\endhead
\bottomrule\noalign{}
\endlastfoot
\textbf{SE-A1} & Lyapunov Descent &
\(\Phi(S_{t+1}, M_{t+1}, f_{\theta_{t+1}}) \leq \Phi(S_t, M_t, f_{\theta_t})\),
with strict inequality unless \(q_{t+1} = q_t\) & SE-1, SE-2 & This is
the \textbf{convergence engine}. Must be proven for the concrete
\(\Phi\) \\
\textbf{SE-A2} & Lyapunov Boundedness &
\(0 \leq \Phi(q) \leq \Phi_0 < \infty\) for all \(q \in \mathcal{Q}\) &
SE-2 & Required for termination bound \\
\textbf{SE-A3} & Finite Data Volume &
\(|\mathcal{D}_{\text{total}}| \leq N_{\text{max}} < \infty\) & SE-2,
SE-3 & Physical constraint \\
\textbf{SE-A4} & Finite Precision & Numerical quantities stored with
precision \(\varepsilon_{\text{mach}} > 0\) & SE-2, SE-3, SE-4, SE-5 &
Computational constraint \\
\textbf{SE-A5} & Lipschitz Gatekeeper & \(w \mapsto S_w\) is
\(L\)-Lipschitz w.r.t. \(\|\cdot\|_\infty\) & SE-4 & Technical condition
for covering number \\
\textbf{SE-A6} & Markovian Evolution &
\(q_{t+1} = G(q_t, \mathcal{D}_{\text{new}, t})\) where
\(\mathcal{D}_{\text{new}, t}\) depends only on \(q_t\) & SE-2 &
Standard in dynamical systems \\
\end{longtable}

\subsubsection{4.3 Assumption Dependency
Graph}\label{assumption-dependency-graph}

\begin{verbatim}
Existing Assumptions:
    A3 (Bounded Loss) ──────────────────────────┐
    A5 (State Homogeneity) ─────────────────────┤
    A1 (Disjoint Training) ──┐                  │
    A2 (Cond. Independence)  ├── (Partial) ─────┤
    A4 (Uniform Noise) ──────┘                  │
    A6 (Balanced Errors) ──── (Partial) ────────┤
                                                │
New Assumptions:                                ▼
    SE-A1 (Lyapunov Descent) ─────────────────→ Theorem SE-1
    SE-A2 (Lyapunov Bounded) ──────────────┐
    SE-A3 (Finite Data) ───────────────────┤
    SE-A4 (Finite Precision) ──────────────┤──→ Theorem SE-2
    SE-A5 (Lipschitz Gatekeeper) ──────────┤
    SE-A6 (Markovian Evolution) ───────────┘
\end{verbatim}

\subsubsection{4.4 Key Observation}\label{key-observation}

The self-evolution theory adds \textbf{6 new assumptions} (SE-A1 through
SE-A6) to the existing 6 assumptions (A1-A6). The new assumptions are of
a \textbf{different character}: - A1-A6 are \textbf{statistical} (about
the data-generating process) - SE-A1/LF are \textbf{dynamical} (about
the evolution algorithm) - SE-A3 and SE-A4 are \textbf{physical} (about
computational resources) - SE-A5 and SE-A6 are \textbf{technical} (for
the covering number and Markovian analysis)

\textbf{Critical gap}: SE-A1 (Lyapunov Descent) is an
\textbf{assumption} in the current theory, not a \textbf{proven
property}. The Lyapunov function \(\Phi\) must be defined and shown to
satisfy the descent property under the self-evolution update rules. This
is the central theoretical challenge for the self-evolution framework.

\begin{center}\rule{0.5\linewidth}{0.5pt}\end{center}

\subsection{5. Proof Gaps Identified}\label{proof-gaps-identified}

\subsubsection{5.1 Gap Classification}\label{gap-classification}

\begin{longtable}[]{@{}
  >{\raggedright\arraybackslash}p{(\linewidth - 6\tabcolsep) * \real{0.2162}}
  >{\raggedright\arraybackslash}p{(\linewidth - 6\tabcolsep) * \real{0.3514}}
  >{\raggedright\arraybackslash}p{(\linewidth - 6\tabcolsep) * \real{0.2703}}
  >{\raggedright\arraybackslash}p{(\linewidth - 6\tabcolsep) * \real{0.1622}}@{}}
\toprule\noalign{}
\begin{minipage}[b]{\linewidth}\raggedright
Gap ID
\end{minipage} & \begin{minipage}[b]{\linewidth}\raggedright
Description
\end{minipage} & \begin{minipage}[b]{\linewidth}\raggedright
Severity
\end{minipage} & \begin{minipage}[b]{\linewidth}\raggedright
Type
\end{minipage} \\
\midrule\noalign{}
\endhead
\bottomrule\noalign{}
\endlastfoot
GAP-1 & Lyapunov function \(\Phi\) not explicitly defined &
\textbf{Critical} & Missing definition \\
GAP-2 & Lyapunov descent property (SE-A1) assumed, not proven &
\textbf{Critical} & Conjectured \\
GAP-3 & Convergence under distribution shift not characterized &
\textbf{High} & Incomplete \\
GAP-4 & Limit cycle conditions not sharply characterized & \textbf{High}
& Incomplete \\
GAP-5 & Lyapunov contraction rate \(\lambda\) not derived &
\textbf{Medium} & Missing bound \\
GAP-6 & \(T^*\) worst-case bound not tight & \textbf{Medium} & Overly
loose \\
GAP-7 & \(\mathcal{S}\) vs \(\hat{\mathcal{S}}_t\) ambiguity &
\textbf{Medium} & Notation fix needed \\
GAP-8 & Precondition checks (Prop 6) not formalized & \textbf{Low} &
Missing procedure \\
GAP-9 & Coupling between \(S_t\) and \(\theta_t\) dynamics &
\textbf{High} & Incomplete \\
GAP-10 & Feature escape mechanism not specified & \textbf{Low} & Future
work \\
\end{longtable}

\subsubsection{5.2 Detailed Gap
Descriptions}\label{detailed-gap-descriptions}

\textbf{GAP-1 (Critical): Lyapunov Function Not Explicitly Defined} The
self-evolution theory asserts the existence of a Lyapunov function
\(\Phi(S_t, M_t, f_{\theta_t})\) but has not provided an explicit
formula. Candidate forms include:

\[\Phi_1 = \sum_{s \in \mathcal{S}} \rho_s \cdot \underbrace{(\mathbb{E}[C(x) \mid x \in s] - \mathbb{E}[S_t(x) \mid x \in s])^2}_{\text{consensus-gatekeeper mismatch}} + \lambda \cdot \underbrace{\text{Err}(f_{\theta_t}, M_t)}_{\text{student prediction error}}\]

\[\Phi_2 = -\text{F1}_{\text{SCX}}(S_t) + \mu \cdot \text{Var}[S_t(x) \mid x \in \mathcal{D}]\]

\[\Phi_3 = \text{KL}(P_{\text{consensus}} \| P_{\text{gatekeeper}}) + \text{Err}(f_{\theta_t}, M_t)\]

Without an explicit \(\Phi\), Theorem SE-1 and SE-2 cannot be rigorously
proven. The descent property (SE-A1) must be derived from the update
rules, not assumed.

\textbf{Status}: Not proven. The theory currently assumes the existence
of a Lyapunov function without constructing one.

\textbf{GAP-2 (Critical): Lyapunov Descent Assumed, Not Proven} Even
given an explicit \(\Phi\), proving \(\Phi(q_{t+1}) \leq \Phi(q_t)\)
requires analyzing the coupled dynamics of: - Gatekeeper \(S_t\) update
(how does \(S_{t+1}\) differ from \(S_t\) given new memory?) - Memory
\(M_t\) update (which samples are added?) - Student \(\theta_t\) update
(how does the NEP retraining change predictions?)

Each of these updates involves stochasticity (new data, random NEP
training). Proving that \(\Phi\) decreases in expectation (or almost
surely) is non-trivial.

\textbf{Status}: Not proven. This is the central technical challenge.

\textbf{GAP-3 (High): Convergence Under Distribution Shift} The NEP
student \(f_{\theta_t}\) explores new configurations over time, changing
the distribution of candidate samples \(\mathcal{C}_t\). This creates a
\textbf{moving target} for the gatekeeper: - \(S_t\) is trained on data
from distribution \(P_t\) (the distribution of \(M_t\)'s samples). - But
it must make decisions on distribution \(P_{t+1}\) (the new candidate
pool). - The distribution shift \(TV(P_t, P_{t+1})\) is not controlled
by the current theory.

This is analogous to the \textbf{off-policy evaluation} problem in
reinforcement learning, or the \textbf{covariate shift} problem in
supervised learning.

\textbf{Status}: Not characterized. The current analysis assumes the
target distribution is fixed.

\textbf{GAP-4 (High): Limit Cycle Conditions Not Characterized} Theorem
SE-2 guarantees convergence to a fixed point, but: - What conditions on
the update dynamics guarantee that the fixed point is \textbf{stable}
(attracting)? - What conditions lead to \textbf{unstable} fixed points
that the system passes through without stopping? - Can the system enter
a near-limit-cycle where \(S_t\) oscillates between several
configurations? (Theorem SE-2 rules out exact cycles via strict Lyapunov
descent, but \(\varepsilon\)-approximate cycles remain possible within
machine precision.)

\textbf{Status}: Partially addressed. Theorem SE-2's finite-state
argument rules out exact cycles, but near-cycles and stability are not
characterized.

\textbf{GAP-5 (Medium): Lyapunov Contraction Rate Not Derived} Theorem
SE-1 (in the architecture files) assumes the Lyapunov function contracts
at rate \(\lambda\):

\[\Phi(q_{t+1}) - \Phi_{\text{opt}} \leq e^{-\lambda}(\Phi(q_t) - \Phi_{\text{opt}})\]

The rate \(\lambda\) is not derived from system parameters. It likely
depends on: - The number of experts \(M\) (more experts \(\to\) faster
contraction) - The state separation gap \(\Delta_{\min}\) (larger gap
\(\to\) faster contraction) - The noise rate \(\eta\) (higher noise
\(\to\) slower contraction) - The validation budget per iteration

\textbf{Status}: Not derived. The contraction rate is a free parameter
in the current theory.

\textbf{GAP-6 (Medium): \(T^*\) Bound Not Tight} The worst-case bound
\(T^* \leq \Phi_0/\varepsilon_{\text{mach}}\) is astronomically loose
(\(\sim 10^{16}\) for double precision). A tighter bound would depend
on: - The number of states \(K_S\) - The number of experts \(M\) - The
per-iteration validation budget \(B\) - The learning rate / gatekeeper
update step size

\textbf{Status}: The theoretical bound is valid but useless for
practical prediction of convergence time.

\textbf{GAP-7 (Medium): \(\mathcal{S}\) vs \(\hat{\mathcal{S}}_t\)
Ambiguity} See Section 3.2. The notation fix is straightforward but
load-bearing for configuration space finiteness.

\textbf{Status}: Needs editorial correction.

\textbf{GAP-8 (Low): Proposition 6 Precondition Not Formalized} The
self-evolution loop should check Proposition 6's stability criterion
before using state-level statistics. The current theory does not
specify: - What threshold of stability triggers the precondition? - What
fallback behavior applies when stability is low? - How does the
stability criterion interact with the Lyapunov function?

\textbf{Status}: Not formalized. This is a practical engineering concern
more than a theoretical gap.

\textbf{GAP-9 (High): Coupled \(S_t\)-\(\theta_t\) Dynamics} The
gatekeeper \(S_t\) and NEP student \(f_{\theta_t}\) co-evolve: 1.
\(S_t\) selects which data to validate. 2. Validated data enters
\(M_t\). 3. \(f_{\theta_t}\) is retrained on \(M_t\). 4. The new
\(f_{\theta_{t+1}}\) explores new configurations. 5. \(S_{t+1}\) must
adapt to the new configuration distribution.

This feedback loop can lead to complex dynamics (oscillations,
bifurcations, chaos) that are not captured by the simple Lyapunov
analysis.

\textbf{Status}: Not characterized. The current analysis treats the
student update as an independent process, ignoring the feedback from
gatekeeper decisions to student training.

\textbf{GAP-10 (Low): Feature Escape Mechanism} Theorem 2's bound is
static: if features are weak, SCX cannot outperform the baseline. The
self-evolution framework should ideally include a mechanism for updating
the feature representation \(\phi\) as the NEP student learns better
internal representations. This is not currently specified.

\textbf{Status}: Not addressed. Acknowledged as future work.

\subsubsection{5.3 Gap Severity Summary}\label{gap-severity-summary}

\begin{verbatim}
GAP-1 ████████████████████░░░░░░░░░░  (Critical, missing definition)
GAP-2 ████████████████████░░░░░░░░░░  (Critical, central technical challenge)
GAP-3 ████████████████░░░░░░░░░░░░░░  (High, under distribution shift)
GAP-4 ████████████████░░░░░░░░░░░░░░  (High, limit cycles)
GAP-9 ████████████████░░░░░░░░░░░░░░  (High, coupled dynamics)
GAP-5 ██████████░░░░░░░░░░░░░░░░░░░░  (Medium, contraction rate)
GAP-6 ██████████░░░░░░░░░░░░░░░░░░░░  (Medium, bound tightness)
GAP-7 ██████████░░░░░░░░░░░░░░░░░░░░  (Medium, notation fix)
GAP-8 ████░░░░░░░░░░░░░░░░░░░░░░░░░░  (Low, procedural)
GAP-10 ████░░░░░░░░░░░░░░░░░░░░░░░░░░ (Low, future work)
\end{verbatim}

\begin{center}\rule{0.5\linewidth}{0.5pt}\end{center}

\subsection{6. Open Problems}\label{open-problems}

The self-evolution theory suggests several open problems for future
investigation:

\subsubsection{\texorpdfstring{Problem 1: Sharp Convergence Rate for
Coupled \(S_t\)-\(\theta_t\)
System}{Problem 1: Sharp Convergence Rate for Coupled S\_t-\textbackslash theta\_t System}}\label{problem-1-sharp-convergence-rate-for-coupled-s_t-theta_t-system}

\textbf{Problem}: Derive the convergence rate of the coupled
gatekeeper-student system \((S_t, f_{\theta_t})\) under the Lyapunov
function \(\Phi\). The rate likely depends on: - The spectral gap of the
state discovery process (from Theorem 5) - The signal-to-noise ratio
\(\Delta_s^2/\sigma^2\) per state - The per-iteration validation budget
\(B\)

\textbf{Conjecture}: For sufficiently separated states
(\(\Delta_{\min}\) large enough), the convergence is linear:

\[\Phi(q_t) - \Phi_{\text{opt}} \leq C \cdot \exp\left(-t \cdot \frac{B \cdot \Delta_{\min}^2}{K_S \cdot \log(1/\varepsilon_{\text{mach}})}\right)\]

\textbf{Difficulty}: Hard. Requires precise control of the state
discovery convergence and the gatekeeper update dynamics.

\subsubsection{Problem 2: Tight Lower Bound on Required Memory
Size}\label{problem-2-tight-lower-bound-on-required-memory-size}

\textbf{Problem}: What is the minimum memory bank size \(|M_t|\)
required to achieve \(\varepsilon\)-optimal gatekeeper performance? The
current bound is trivial (\(N_{\text{max}}\)). A tight bound would
depend on: - The number of states \(K_S\) - The within-state variance of
expert predictions - The noise rate \(\eta\)

\textbf{Conjecture}: The required memory scales as:

\[|M_{\min}| = \Omega\left(K_S \cdot \frac{1}{\varepsilon^2} \cdot \log\frac{K_S}{\delta}\right)\]

where \(\varepsilon\) is the target gatekeeper accuracy and \(\delta\)
is the failure probability.

\textbf{Difficulty}: Moderate. Related to sample complexity in
state-level estimation.

\subsubsection{Problem 3: Phase Transition Between Four Convergence
Regimes}\label{problem-3-phase-transition-between-four-convergence-regimes}

\textbf{Problem}: Characterize the conditions under which the
self-evolution exhibits each of four regimes:

\begin{enumerate}
\def\labelenumi{\arabic{enumi}.}
\tightlist
\item
  \textbf{Fast convergence}: Gatekeeper improves rapidly, student
  discovers new states quickly, fixed point reached in \(O(K_S)\)
  iterations.
\item
  \textbf{Slow convergence}: Gatekeeper improves slowly, student
  discovers new states incrementally, fixed point reached in
  \(O(K_S/\varepsilon_{\text{mach}})\) iterations.
\item
  \textbf{Limit cycle}: Gatekeeper oscillates between configurations
  (near-cycle within machine precision). \(\Phi\) decreases overall but
  individual steps may increase.
\item
  \textbf{Stall}: Gatekeeper cannot improve because the consensus signal
  is too weak (Theorem 1's \(\Delta_s \approx 0\) for all \(s\)). The
  system reaches a false fixed point that is far from optimal.
\end{enumerate}

\textbf{Difficulty}: Hard. Requires a dynamical systems analysis of the
SCX update rules.

\subsubsection{Problem 4: Optimal Gatekeeper Update
Frequency}\label{problem-4-optimal-gatekeeper-update-frequency}

\textbf{Problem}: How often should the gatekeeper be updated relative to
the student? Currently, the framework updates both every iteration.
Alternatives include: - \textbf{Slow gatekeeper, fast student}: Update
the student every iteration, gatekeeper every \(K\) iterations. -
\textbf{Fast gatekeeper, slow student}: Update the gatekeeper after
every validation event, student only when memory bank has grown
significantly. - \textbf{Asynchronous updates}: Update gatekeeper and
student independently based on information-theoretic triggers.

\textbf{Conjecture}: The optimal update frequency depends on the
relative rates of information gain. When the consensus signal is
changing rapidly (early in evolution), update the gatekeeper frequently.
When the consensus signal is stable, update the student more frequently.

\textbf{Difficulty}: Moderate. An empirical study could establish
guidelines even without a fully rigorous theory.

\subsubsection{Problem 5: Finite-Time Guarantees (Non-Asymptotic
Bounds)}\label{problem-5-finite-time-guarantees-non-asymptotic-bounds}

\textbf{Problem}: Provide non-asymptotic bounds on the Lyapunov gap
\(\Phi(q_t) - \Phi_{\text{opt}}\) after \(t\) iterations. The current
theory provides only asymptotic convergence (Theorem SE-1) and
worst-case finite-time termination (Theorem SE-2). Neither gives a
useful bound for practical \(t\) (e.g., \(t \leq 1000\)).

\textbf{Desired result}: For any \(t < T^*\), with probability at least
\(1-\delta\):

\[\Phi(q_t) - \Phi_{\text{opt}} \leq \Phi(q_0) \cdot \exp\left(-c \cdot \frac{t \cdot B \cdot \Delta_{\min}^2}{K_S \cdot d_\phi}\right) + \varepsilon_{\text{stat}}\]

where \(\varepsilon_{\text{stat}}\) is the statistical error due to
finite samples in \(M_t\).

\textbf{Difficulty}: Hard. Requires concentration bounds for a
sequential, adaptive process.

\begin{center}\rule{0.5\linewidth}{0.5pt}\end{center}

\subsection{7. Numerical Verification
Suggestions}\label{numerical-verification-suggestions}

\subsubsection{7.1 Experiment 1: Lyapunov Descent
Verification}\label{experiment-1-lyapunov-descent-verification}

\textbf{Objective}: Verify that the Lyapunov function \(\Phi\) (once
defined) decreases monotonically.

\textbf{Setup}: - Synthetic data with known states, known noise rates
\(\eta\), and known expert error rates \(\varepsilon_m\). - Run the
self-evolution loop for \(T = 100\) iterations. - Track
\(\Phi(S_t, M_t, f_{\theta_t})\) at each iteration.

\textbf{Metrics}: - Monotonicity: \(\mathbb{P}(\Phi_{t+1} < \Phi_t)\)
across iterations. - Descent rate:
\(\lambda_t = \log(\Phi_t / \Phi_{t-1})\) per iteration. - Plateau
detection: identify \(T^*\) where \(\Phi\) stabilizes (change
\(< \varepsilon_{\text{mach}}\) for 5 consecutive iterations).

\textbf{Expected result}: \(\Phi\) decreases monotonically (but may
plateau before reaching the global minimum). The descent rate depends on
\(\Delta_{\min}\) and the validation budget.

\subsubsection{7.2 Experiment 2: Theorem SE-2 Termination
Bound}\label{experiment-2-theorem-se-2-termination-bound}

\textbf{Objective}: Test whether the system reaches a fixed point in
finite time, and measure \(T^*\) as a function of problem parameters.

\textbf{Setup}: - Vary: \(K_S\) (2 to 20), \(\eta\) (0.01 to 0.4),
\(\Delta_{\min}\) (0.1 to 0.5), \(M\) (5 to 50). - For each
configuration, run the self-evolution loop until \(\Phi\) stabilizes.

\textbf{Metrics}: - \(T^*\) as a function of \(K_S\), \(\eta\),
\(\Delta_{\min}\), \(M\). - Compare to the theoretical bound
\(T^* \leq \Phi_0/\varepsilon_{\text{mach}}\).

\textbf{Expected result}: \(T^*\) scales like \(O(K_S/\Delta_{\min}^2)\)
for large \(\Delta_{\min}\), transitioning to exponential scaling as
\(\Delta_{\min} \to 0\).

\subsubsection{7.3 Experiment 3: Fixed Point Quality vs.~Theorem
4'}\label{experiment-3-fixed-point-quality-vs.-theorem-4}

\textbf{Objective}: Verify that the fixed point's noise detection F1
matches Theorem 4'\,'s optimal constant.

\textbf{Setup}: - Same synthetic data as Experiment 1. - At the fixed
point \(T^*\), compute the empirical F1 of \(S_{T^*}(x)\) as a noise
detector. - Compare to Theorem 4'\,'s prediction:
\(1 - \text{F1} \approx C_{\min}/(\eta \cdot e^{M\kappa}\sqrt{2\pi M})\).

\textbf{Metrics}: - Ratio: empirical F1 gap / theoretical F1 gap. Should
approach 1 as \(M \to \infty\). - Sensitivity to \(p_0\), \(p_1\),
\(\eta\) (the three parameters of Theorem 4').

\textbf{Expected result}: Agreement within statistical error for
\(M \geq 20\). Divergence for small \(M\) where asymptotic
approximations break down.

\subsubsection{7.4 Experiment 4: Feature Bottleneck (Theorem 2
Confirmation)}\label{experiment-4-feature-bottleneck-theorem-2-confirmation}

\textbf{Objective}: Verify that the self-evolution converges to a fixed
point that respects Theorem 2's bound.

\textbf{Setup}: - Generate features \(\phi\) with controlled mutual
information \(\delta = I(\phi; S)\). - Run self-evolution for features
with \(\delta \in \{0.001, 0.01, 0.1, 0.5, 1.0\}\) nats. - Compare the
fixed-point F1 to Theorem 2's bound:
\(F1_{\text{base}} + C_F\sqrt{\delta/2}\).

\textbf{Metrics}: - Does fixed-point F1 exceed the bound? (Should not,
if theory is correct.) - How close does the fixed point come to the
bound for large \(\delta\)? - Does the self-evolution detect the feature
bottleneck (e.g., through Proposition 6's stability score)?

\textbf{Expected result}: The fixed-point F1 remains below the Theorem 2
bound. For small \(\delta\), the fixed point F1 is close to the bound
(the gatekeeper is not the bottleneck; features are). For large
\(\delta\), the fixed point may be significantly below the bound (the
gatekeeper or student is the bottleneck).

\subsubsection{7.5 Experiment 5: External Validation Escape from
Incompleteness}\label{experiment-5-external-validation-escape-from-incompleteness}

\textbf{Objective}: Test the Claim SE-C2 about external validation
escaping incompleteness.

\textbf{Setup}: - Run self-evolution to fixed point \(T^*\) without
external validation. - Inject one round of external NEP validation
(simulated: reveal true labels for a small set of samples). - Continue
self-evolution. Does it escape the fixed point?

\textbf{Metrics}: - Escape success rate: fraction of fixed points broken
by external validation. - Improvement after escape:
\(\Phi_{\text{after}} - \Phi_{\text{at fixed point}}\). - Sensitivity to
validation budget (how many externally validated samples are needed to
escape).

\textbf{Expected result}: A small amount of external validation (1-5\%
of the total data) is sufficient to break most suboptimal fixed points.
The required budget depends on the noise rate \(\eta\) and the state
separation \(\Delta_{\min}\).

\begin{center}\rule{0.5\linewidth}{0.5pt}\end{center}

\subsection{8. Honest Assessment}\label{honest-assessment}

\subsubsection{8.1 Classification Scheme}\label{classification-scheme}

Each claim in the self-evolution theory is classified into one of three
categories:

\begin{longtable}[]{@{}
  >{\raggedright\arraybackslash}p{(\linewidth - 4\tabcolsep) * \real{0.3846}}
  >{\raggedright\arraybackslash}p{(\linewidth - 4\tabcolsep) * \real{0.2692}}
  >{\raggedright\arraybackslash}p{(\linewidth - 4\tabcolsep) * \real{0.3462}}@{}}
\toprule\noalign{}
\begin{minipage}[b]{\linewidth}\raggedright
Category
\end{minipage} & \begin{minipage}[b]{\linewidth}\raggedright
Label
\end{minipage} & \begin{minipage}[b]{\linewidth}\raggedright
Meaning
\end{minipage} \\
\midrule\noalign{}
\endhead
\bottomrule\noalign{}
\endlastfoot
\textbf{Rigorous Theory} & \(\checkmark\) & Fully proven with explicit
assumptions and mathematical derivation. No gaps. \\
\textbf{Formal Conjecture} & \(\triangle\) & Mathematically well-posed
claim with plausible reasoning, but proof steps are missing or
assumptions are not fully checked. \\
\textbf{Plausible Hypothesis} & \(\circ\) & Well-motivated claim based
on intuition or analogy. Formalization is incomplete or at an early
stage. \\
\end{longtable}

\subsubsection{8.2 Claim-by-Claim
Assessment}\label{claim-by-claim-assessment}

\begin{longtable}[]{@{}
  >{\raggedright\arraybackslash}p{(\linewidth - 6\tabcolsep) * \real{0.1944}}
  >{\raggedright\arraybackslash}p{(\linewidth - 6\tabcolsep) * \real{0.2778}}
  >{\raggedright\arraybackslash}p{(\linewidth - 6\tabcolsep) * \real{0.2778}}
  >{\raggedright\arraybackslash}p{(\linewidth - 6\tabcolsep) * \real{0.2500}}@{}}
\toprule\noalign{}
\begin{minipage}[b]{\linewidth}\raggedright
Claim
\end{minipage} & \begin{minipage}[b]{\linewidth}\raggedright
Category
\end{minipage} & \begin{minipage}[b]{\linewidth}\raggedright
Evidence
\end{minipage} & \begin{minipage}[b]{\linewidth}\raggedright
Key Gap
\end{minipage} \\
\midrule\noalign{}
\endhead
\bottomrule\noalign{}
\endlastfoot
\textbf{Theorem SE-1} (Lyapunov Convergence) & \(\triangle\) Formal
Conjecture & The Lyapunov descent idea is standard in optimization;
adaptation to SCX is plausible. & \(\Phi\) not explicitly defined;
descent property not proven for coupled dynamics. \\
\textbf{Theorem SE-2} (Completeness Bound) & \(\checkmark\) Rigorous
Theory & Finite configuration space + finite-state dynamics + strict
Lyapunov descent \(\implies\) finite-time fixed point. The logic is
sound. & Depends on SE-A1 (Lyapunov descent) being satisfied. Without
it, SE-2's conclusion weakens to ``convergence to a cycle.'' \\
\textbf{Proposition SE-3} (Finite Configuration Space) & \(\checkmark\)
Rigorous Theory & Trivially true under finite data, finite precision,
finite parameterization. No gaps. & None. \\
\textbf{Proposition SE-4} (Covering Number) & \(\checkmark\) Rigorous
Theory & Standard result from empirical process theory; applies directly
to parametric function classes. & The Lipschitz constant \(L\) may be
problem-dependent. \\
\textbf{Proposition SE-5} (Metric Entropy) & \(\checkmark\) Rigorous
Theory & Direct corollary of SE-4. & None. \\
\textbf{Claim SE-C1} (True Statements Unprovable Within System) &
\(\circ\) Plausible Hypothesis & The analogy to Godel incompleteness is
structurally suggestive but not formally established. & The SCX system
is not a formal logic; ``provability'' is not rigorously defined. \\
\textbf{Claim SE-C2} (Cannot Self-Certify) & \(\circ\) Plausible
Hypothesis & Self-referential evaluation is epistemically limited. The
intuition is sound. & ``Self-certification'' is not formally defined.
The claim is more philosophical than mathematical. \\
\textbf{External Validation Escape} & \(\circ\) Plausible Hypothesis &
Introducing new information can break fixed points. This is intuitively
true. & Formal mechanism for escaping fixed points via external
validation is not specified. \\
\textbf{Finite \(T^*\) Bound} & \(\checkmark\) Rigorous Theory &
\(T^* \leq \Phi_0/\varepsilon_{\text{mach}}\) is a straightforward
consequence of finite-state dynamics. & The bound is astronomically
loose. \\
\textbf{Convergence Regimes} (fast/slow/cycle/stall) & \(\circ\)
Plausible Hypothesis & The four regimes are qualitatively plausible. &
No formal characterization of regime boundaries. \\
\textbf{\(S_t\)-\(\theta_t\) Coupling} & \(\triangle\) Formal Conjecture
& The feedback loop exists and matters. Analyzing it requires dynamical
systems theory. & No formal analysis has been attempted. \\
\end{longtable}

\subsubsection{8.3 Overall Assessment}\label{overall-assessment}

\begin{verbatim}
Rigorous Theory    ████████░░░░░░████  (~40% of claims)
Formal Conjecture  ██████░░░░░░░░░░░░  (~20% of claims)
Plausible Hypoth.  ████████░░░░░░░░░░  (~40% of claims)
\end{verbatim}

The self-evolution theory has a solid core of rigorous results: -
\textbf{Proposition SE-3}: Finite configuration space (trivially true) -
\textbf{Proposition SE-4/5}: Covering number and metric entropy
(standard math) - \textbf{Theorem SE-2}: Finite-time termination (sound
logic, but conditional on SE-A1)

The theory's main weakness is the \textbf{central dependence on an
unproven Lyapunov descent assumption}: - Theorem SE-1 assumes rather
than proves monotonic convergence. - Theorem SE-2's strongest form
depends on SE-A1 being satisfied. - Without SE-A1, SE-2 only guarantees
convergence to a cycle (still finite-time, but weaker).

The Godel incompleteness analogy (Claims SE-C1, SE-C2) is
\textbf{illuminating but not rigorous}. It serves as a conceptual guide
rather than a mathematical result.

\subsubsection{8.4 Priority
Recommendations}\label{priority-recommendations}

To move from the current conjectural state to a rigorous theory, the
following steps are recommended in priority order:

\begin{longtable}[]{@{}lll@{}}
\toprule\noalign{}
Priority & Action & Affected Claims \\
\midrule\noalign{}
\endhead
\bottomrule\noalign{}
\endlastfoot
\textbf{P0} & Define \(\Phi\) explicitly & SE-1, SE-2 \\
\textbf{P1} & Prove Lyapunov descent for concrete \(\Phi\) & SE-1 \\
\textbf{P2} & Analyze \(S_t\)-\(\theta_t\) coupling & SE-1 (tightness),
GAP-9 \\
\textbf{P3} & Derive contraction rate \(\lambda\) & SE-1 (practical
bound) \\
\textbf{P4} & Formalize external validation escape & SE-C2 \\
\textbf{P5} & Tighten \(T^*\) bound & SE-2 (practical utility) \\
\end{longtable}

\begin{center}\rule{0.5\linewidth}{0.5pt}\end{center}

\subsection{9. Summary of Verification
Status}\label{summary-of-verification-status}

\subsubsection{9.1 Pass / Fail / Gap
Summary}\label{pass-fail-gap-summary}

\begin{longtable}[]{@{}
  >{\raggedright\arraybackslash}p{(\linewidth - 4\tabcolsep) * \real{0.4231}}
  >{\raggedright\arraybackslash}p{(\linewidth - 4\tabcolsep) * \real{0.3077}}
  >{\raggedright\arraybackslash}p{(\linewidth - 4\tabcolsep) * \real{0.2692}}@{}}
\toprule\noalign{}
\begin{minipage}[b]{\linewidth}\raggedright
Component
\end{minipage} & \begin{minipage}[b]{\linewidth}\raggedright
Status
\end{minipage} & \begin{minipage}[b]{\linewidth}\raggedright
Notes
\end{minipage} \\
\midrule\noalign{}
\endhead
\bottomrule\noalign{}
\endlastfoot
New definitions (Files 1-6) & \textbf{PASS} (with minor warnings) &
\(\mathcal{S}\) vs \(\hat{\mathcal{S}}_t\) needs fix \\
Theorem SE-1 (Lyapunov) & \textbf{GAP} (missing \(\Phi\) definition,
descent unproven) & Central open problem \\
Theorem SE-2 (Completeness) & \textbf{PASS} (conditional on SE-A1) &
Valid logic, loose bound \\
Prop SE-3 (Finite Space) & \textbf{PASS} & Trivially correct \\
Prop SE-4 (Covering Number) & \textbf{PASS} & Standard result \\
Prop SE-5 (Metric Entropy) & \textbf{PASS} & Direct corollary \\
Claims SE-C1/C2 (Godel) & \textbf{CONJECTURAL} & Illuminating but not
formal \\
Cross-reference with Thm 1 & \textbf{PASS} (dependency should be
explicit) & -- \\
Cross-reference with Thm 2 & \textbf{PASS} (limitation respected) &
-- \\
Cross-reference with Thm 3 & \textbf{PASS} (ambiguity persists) & -- \\
Cross-reference with Thm 4' & \textbf{PASS} (connection needs explicit
tie) & -- \\
Cross-reference with Thm 5 & \textbf{PASS} (precondition recommended) &
-- \\
Notation consistency & \textbf{WARNING} (see conflicts) &
\(\mathcal{S}\), \(\theta\) conflicts \\
Assumption completeness & \textbf{ADEQUATE} & 12 total assumptions (6
old + 6 new) \\
Gap identification & \textbf{COMPLETE} & 10 gaps identified (2 critical,
3 high, 2 medium, 3 low) \\
Open problems & \textbf{FORMULATED} & 5 problems defined \\
Numerical suggestions & \textbf{PROVIDED} & 5 experiments described \\
\end{longtable}

\subsubsection{9.2 Final Assessment}\label{final-assessment}

The SCX self-evolution theory is at a \textbf{preliminary but promising
stage}:

\begin{itemize}
\item
  \textbf{Strengths}: The termination guarantee (Theorem SE-2) is
  rigorous and provides a principled foundation. The cross-references to
  existing theorems are consistent. The conceptual connections to
  AlphaZero, BO, AL, and Solomonoff induction are well-drawn.
\item
  \textbf{Weaknesses}: The central convergence claim (Theorem SE-1) is
  not yet proven --- the Lyapunov function remains undefined and its
  descent property unproven. The coupled dynamics of gatekeeper and
  student are not analyzed. Several claims are plausible hypotheses
  rather than theorems.
\item
  \textbf{Recommendation}: The theory is appropriate for a pre-print or
  workshop paper with honest labeling of which parts are rigorous and
  which are conjectural. For a journal submission, the Lyapunov gap
  (GAP-1, GAP-2) must be closed, and the \(S_t\)-\(\theta_t\) coupling
  (GAP-9) must be at least partially characterized.
\end{itemize}

\begin{center}\rule{0.5\linewidth}{0.5pt}\end{center}

\emph{End of 09\_verification\_report.md}

\end{document}
